% !PRINTED 634
\weberpage{634}
Gehen wir nun zur Substitution $t$ über und setzen in
\weq{2}{(\varepsilon, x) = x_0 + \sum_{1,\,n-1}^{h}\varepsilon^h x_h}
$x_{gh}$ für $x_h$, so geht diese Function über in
\weq{3}{x_0 + \sum_{1,\,n}^{h}\varepsilon^h x_{gh}.} % [sic] Print has lower limits $1,n$ although the residue system excludes $0$; preserved.
In dieser Summe kann $h$ ein beliebiges Restsystem nach dem Modul $n$ durchlaufen mit Ausschluss von $0$, und wir können also $h$ durch $g^{-1}h$ ersetzen; dadurch wird die Summe (3)
\[
x_0 + \sum^{h}\varepsilon^{g^{-1}h}x_h = (\varepsilon^{g^{-1}}, x).
\]
Wendet man die Permutationen $s$ und $t$ wiederholt an, so ergiebt sich der Einfluss von $s^\lambda, t^\lambda$, der in den Vertauschungen
\[
(\varepsilon, x), \varepsilon^{-\lambda}(\varepsilon, x)
\]
\[
(\varepsilon, x), (\varepsilon^{g^{-\lambda}}, x)
\]
besteht, speciell also für $\lambda = -1$
\[
(\varepsilon, x), \varepsilon(\varepsilon, x)
\]
\[
(\varepsilon, x), (\varepsilon^{g}, x),
\]
und wir bekommen den folgenden Satz:

\gesperrt{4. Die erzeugenden Permutationen der metacyklischen Gruppe, $s, t$, bewirken die Vertauschungen
\[
(s)\colon (\varepsilon, x), \varepsilon^{-1}(\varepsilon, x); \qquad (s^{-1})\colon (\varepsilon, x), \varepsilon(\varepsilon, x)
\]
\[
(t)\colon (\varepsilon, x), (\varepsilon^{g^{-1}}, x); \qquad (t^{-1})\colon (\varepsilon, x), (\varepsilon^{g}, x).
\]}

Aus 2. und 4. aber ergiebt sich, weil $s$ die erzeugende Permutation der cyklischen Gruppe ist, die Folgerung:

\gesperrt{5. Stellt man die Functionen
\weq{4}{(\varepsilon, x)^n = \Phi(\varepsilon), \quad (\varepsilon^\lambda, x)(\varepsilon, x)^{-\lambda} = F(\varepsilon)}
in der Normalform [\S.\,183, (3)] dar, so sind die Coëfficienten cyklische Functionen von $x_0, x_1 \ldots x_{n-1}$ (\S.\,163).}

Hierin kann $\lambda$ jede beliebige ganze Zahl, die nicht durch $n$ theilbar ist, bedeuten.

Aus der Function $F(\varepsilon)$ entspringen, wenn wir $\lambda = g$ annehmen und für $\varepsilon$ seine $n-1$ verschiedenen Werthe setzen, $n-1$ Functionen, die alle durch $s$ ungeändert bleiben, und die wir folgendermaassen bezeichnen wollen:
%CONT>>
